% =============================================================================
% A Reproducible Pipeline for Architecture Comparison in Low-Resource
% Language Modelling
% Class final project, Natural Language Processing, Insper, 2026
% =============================================================================

\documentclass[11pt]{article}

% --- packages -----------------------------------------------------------------
\usepackage[a4paper, margin=1in]{geometry}
\usepackage[utf8]{inputenc}
\usepackage[T1]{fontenc}
\usepackage{mathptmx}                    % Times-like body + math (legacy, works everywhere)
\usepackage{microtype}                    % nicer typography
\usepackage{graphicx}
\usepackage{booktabs}
\usepackage{caption}
\usepackage{subcaption}
\usepackage{amsmath}                       % \text in math, etc.
\usepackage[hidelinks]{hyperref}
\usepackage{xcolor}
\usepackage{titlesec}
\usepackage{titling}
\usepackage{enumitem}
\usepackage{listings}
\usepackage{authblk}
\usepackage{abstract}

% --- styling: clean, single column, Nature-ish -------------------------------
\setlength{\droptitle}{-1.2cm}
\pretitle{\begin{flushleft}\large\bfseries}
\posttitle{\par\end{flushleft}\vskip 0.4em}
\preauthor{\begin{flushleft}\normalsize}
\postauthor{\par\end{flushleft}}
\predate{\begin{flushleft}\small\itshape}
\postdate{\par\end{flushleft}\vskip 0.6em}

\titleformat*{\section}{\large\bfseries}
\titleformat*{\subsection}{\normalsize\bfseries}
\titleformat*{\subsubsection}{\normalsize\bfseries\itshape}

% Tighter spacing around section heads
\titlespacing*{\section}{0pt}{1.2em}{0.4em}
\titlespacing*{\subsection}{0pt}{0.9em}{0.3em}

% Abstract: bold heading, no italic, slightly indented body
\renewcommand{\abstractnamefont}{\normalsize\bfseries}
\renewcommand{\abstracttextfont}{\normalsize}
\setlength{\absleftindent}{0pt}
\setlength{\absrightindent}{0pt}

\setlength{\parskip}{0.4em}
\setlength{\parindent}{0pt}

% Code listings
\lstset{
  basicstyle=\ttfamily\small,
  breaklines=true,
  frame=single,
  framerule=0.4pt,
  rulecolor=\color{black!40},
  xleftmargin=0.5em,
  xrightmargin=0.5em,
  aboveskip=0.6em,
  belowskip=0.6em,
}

% =============================================================================
\title{A Reproducible Pipeline for Architecture Comparison in
       Low-Resource Language Modelling:\\
       A Case Study on xLSTM versus Transformer}

\author[1]{João~Pedro~Rodrigues~dos~Santos}
\affil[1]{Computer Engineering, Insper, São Paulo, Brazil}
\date{May 2026}

% =============================================================================
\begin{document}

\maketitle

\begin{abstract}
\noindent
Fair comparison of neural network architectures is difficult: differences in
tokenizers, hyperparameters, data preparation, and infrastructure routinely
conflate with architectural differences in published comparisons. This paper
presents an open-source pipeline designed to isolate architecture as the
sole variable in language modelling experiments at small scales
($1$--$10$\,M parameters, $1.5$--$130$\,M training tokens). The pipeline
standardises raw-text ingestion, byte-level BPE tokenizer training, model
construction, training, and deterministic perplexity evaluation across
user-supplied architectures. We illustrate the pipeline by comparing
xLSTM~\cite{beck2024} and a GPT-style Transformer at five parameter scales
across two hardware regimes (CPU and GPU), using nineteen 19th-century
Portuguese novels from Project Gutenberg as the training corpus, and
complement perplexity with three downstream linear-probe tasks: author
classification (in-domain), sentiment classification on a
MarianMT~\cite{junczys2018marian,tiedemann2020opus} machine-translated
subset of IMDB~\cite{maas2011imdb}, and review-score classification on
native-Portuguese Olist e-commerce data. We observe a scale-dependent
crossover in held-out perplexity: the Transformer wins at the smallest
scale, the xLSTM wins at the CPU medium scale, the Transformer wins again
at $10$\,M parameters on GPU, where it however overfits less aggressively
than the xLSTM. On the downstream probes, the xLSTM leads on author
classification and on translated IMDB sentiment, while the two architectures
are at statistical parity on native-Portuguese Olist sentiment. We
emphasise a methodological caveat that is rarely controlled in published
comparisons: training throughput between the two architectures in our
setting differs by $\sim$$5$--$8\times$ on CPU and by $\sim$$45$--$80\times$
on GPU due to implementation choices, not architectural cost. The pipeline
is released as open-source software, ready for extension to additional
architectures, corpora, and scales.
\end{abstract}

% =============================================================================
\section{Introduction}

The dominant architecture in modern language modelling is the
Transformer~\cite{vaswani2017}, which underpins both encoder
models such as BERT~\cite{devlin2019} and the GPT family of decoders.
Despite its dominance, recurrent architectures have not disappeared. The
xLSTM family, introduced recently by Beck et al.~\cite{beck2024}, revisits
the LSTM with two new components: the \emph{sLSTM}, which augments the
classical LSTM gating with exponential gates and a stabilisation mechanism;
and the \emph{mLSTM}, which replaces the scalar memory cell with a matrix
memory and is fully parallelisable. The authors argue that xLSTMs achieve
parameter efficiency competitive with Transformers in language modelling,
particularly at scale.

Most published comparisons between recurrent and attention-based language
models target the regime in which both architectures shine: large models
trained on large corpora. Considerably less is known about the
\emph{extremely low-resource regime}, where models contain a few million
parameters and are trained on a few million tokens. This regime is relevant
in two practical settings: pedagogical and prototyping experiments where
compute budgets are tight, and bootstrapping in low-resource languages
where pretraining data is fundamentally limited.

A second, methodological problem is that comparisons published across
papers are difficult to interpret because they typically differ in
tokenizer, hyperparameters, data preparation, evaluation set, and software
stack. Conclusions about architectural superiority routinely confound with
these other choices. The community has noted this reproducibility gap, but
public, end-to-end pipelines that explicitly isolate architecture as
the sole variable remain scarce.

\paragraph{Contribution.} The primary contribution of this work is
methodological: an open-source, end-to-end pipeline that isolates
architecture as the only variable in small-scale language modelling
experiments. The pipeline standardises every other component:
\begin{itemize}[leftmargin=1.2em, itemsep=0pt, topsep=0pt]
  \item raw-text ingestion and cleaning (Project Gutenberg boilerplate
        stripping);
  \item byte-level BPE tokenizer training;
  \item model factory parameterised by a single \texttt{arch} flag;
  \item training loop (AdamW, cosine schedule, gradient clipping,
        identical seeds);
  \item deterministic perplexity evaluation on a held-out corpus;
  \item logging in a single \texttt{log.jsonl} format suitable for
        cross-run aggregation and plotting.
\end{itemize}
As a secondary contribution, we apply the pipeline to xLSTM and a
GPT-style Transformer at five parameter scales across two hardware
regimes (CPU: $\sim$$1.5$\,M and $\sim$$4.5$\,M; GPU: $\sim$$2.5$\,M,
$\sim$$6.5$\,M, and $\sim$$10$\,M), pairing held-out perplexity with
three linear-probe downstream tasks (author classification, sentiment
on machine-translated IMDB-PT, review-score on native-Portuguese Olist
e-commerce data), and report empirical observations. These
observations are intentionally framed as an illustration of the
pipeline rather than as definitive claims about architectural
superiority.

% =============================================================================
\section{Related Work}

\paragraph{Transformer-based language modelling.}
The Transformer~\cite{vaswani2017} is currently the default architecture
for both encoder-style models~\cite{devlin2019} and autoregressive
decoders, including the GPT line~\cite{radford2019}. The decoder
configuration we adopt as a baseline (pre-norm layers, multi-head causal
attention, GELU MLPs) is essentially the small-scale GPT recipe popularised
by Karpathy's \texttt{nanoGPT}~\cite{nanogpt}, which we use as a reference
implementation.

\paragraph{xLSTM.}
xLSTM~\cite{beck2024} comprises two block types. The sLSTM block uses
exponential gating to address the limited storage capacity of classical
LSTMs and is recurrent in nature. The mLSTM block stores associations in
a matrix, exposes a parallel formulation analogous to attention, and
provides parameter-efficient long-range memory. The original work
demonstrates competitive performance with Transformers on standard
language modelling benchmarks at scale.

\paragraph{Tokenization and tokenizer-induced confounds.}
Sennrich et al.~\cite{sennrich2016} introduced byte-pair encoding for
neural machine translation; the byte-level variant we employ is standard
in modern decoders~\cite{radford2019}. We note that comparisons across
architectures often differ in tokenizer choice, which is a known confound
since tokenizer determines effective vocabulary, sequence length, and the
intrinsic difficulty of the language modelling objective.

% =============================================================================
\section{Methodology: A Reproducible Comparison Pipeline}

The pipeline is organised in four stages
(Figure~\ref{fig:pipeline}). Stage 1 ingests raw text from a list of
Project Gutenberg book identifiers and strips the standard boilerplate
markers. Stage 2 trains a byte-level BPE tokenizer with a configurable
vocabulary size on the materialised corpus. Stage 3 constructs the model
through an architecture-agnostic factory that dispatches on a single
\texttt{arch} field in the training configuration; this is the only
configuration field that differs between comparable runs. Stage 4 runs the
training loop and produces a single \texttt{log.jsonl} per run that is
consumed downstream by analysis scripts.

\begin{figure}[h]
\centering
\fbox{\parbox{0.95\linewidth}{\centering\small
\textsf{raw books}~$\rightarrow$~\textsf{cleaning}~$\rightarrow$~\textsf{tokenizer training}~$\rightarrow$~\textsf{model factory (arch=xlstm \textbar{} transformer)}~$\rightarrow$~\textsf{training loop}~$\rightarrow$~\textsf{log.jsonl}~$\rightarrow$~\textsf{aggregated comparison}
}}
\caption{The four-stage comparison pipeline. The dispatch on
  \texttt{arch} is the only locus where architectures diverge; every
  other component is shared.}
\label{fig:pipeline}
\end{figure}

\subsection{Data}

The training corpus consists of nineteen novels from Project Gutenberg,
totalling \textasciitilde 7.5\,MB of cleaned UTF-8 text
(\textasciitilde 1.9\,M tokens at the byte-level BPE described below).
The selection covers six authors across multiple literary movements of
19th-century Brazilian and Portuguese literature: \emph{Dom Casmurro},
\emph{Quincas Borba}, \emph{Helena}, \emph{A Mão e a Luva},
\emph{Esaú e Jacó} and \emph{Memorial de Aires} by Machado de Assis
(Brazilian realism); \emph{O Primo Basílio}, \emph{O Mandarim},
\emph{Os Maias}, \emph{A Relíquia} and \emph{A Cidade e as Serras} by
Eça de Queirós (Portuguese realism); \emph{Iracema}, \emph{O Guarany}
(Vols. 1 and 2), \emph{Ubirajara} and \emph{A Pata da Gazela} by José
de Alencar (Brazilian Romantic indianism); \emph{Amor de Perdição} by
Camilo Castelo Branco (Portuguese romanticism); \emph{A Escrava Isaura}
by Bernardo Guimarães (Brazilian abolitionist romanticism); and \emph{O
Ateneu} by Raul Pompéia (Brazilian late-19th-century prose). The
held-out evaluation corpus is \emph{Memórias Póstumas de Brás Cubas}
by Machado de Assis. The holdout shares an author with six training
works (the dominant author in the period) but is itself work-disjoint;
the literature available in the public domain in Portuguese imposes
limits on the achievable holdout disjointness.

The Project Gutenberg standard boilerplate (English-language license
header and footer) is stripped using regular expressions matching the
canonical \texttt{*** START} and \texttt{*** END} markers. Without this
step, models learn to predict English legal text at the boundaries.

\subsection{Tokenizer}

We train a single byte-level BPE tokenizer with vocabulary size 4{,}000
on the same corpus that will be used for training, using the
\texttt{tokenizers} library~\cite{huggingfacetokenizers}. The byte-level
variant is chosen for two reasons: (i)~it eliminates the possibility of
\texttt{<unk>} tokens, which is desirable for accented Portuguese
orthography; and (ii)~it permits a uniform treatment across the (eventual)
multilingual extension of the pipeline. The vocabulary size of 4{,}000
is small by modern standards but appropriate for the corpus size and
deliberately stresses the language modelling objective.

\subsection{Architectures}

Two architectures are compared. Both share the same input vocabulary
size, embedding dimension, context length, and output head; they differ
only in their core block.

\paragraph{xLSTM.}
We use the official \texttt{xlstm} package~\cite{beck2024}. Each model
contains one sLSTM block and the remainder mLSTM blocks. The sLSTM block
is placed at the second position of the stack, following the recipe in
the original paper. We use the \texttt{vanilla} backend (pure PyTorch),
since the high-performance \texttt{cuda} backend requires CUDA Compute
Capability \(\ge\)\,8.0 and our experiments are conducted on CPU. This
choice is methodologically conservative for the speed comparison
discussed in Section~\ref{sec:throughput}.

\paragraph{Transformer.}
We implement a GPT-style decoder in pure PyTorch: token and learned
positional embeddings, a stack of pre-norm blocks each containing
multi-head causal self-attention (computed via
\texttt{F.scaled\_dot\_product\_attention}) and a GELU MLP, a final
layer norm, and an untied output projection. Weight initialisation
follows GPT-2~\cite{radford2019}: \(\mathcal{N}(0, 0.02^2)\) for weights and
zero for biases.

\paragraph{Parameter pairing.}
At each scale, the two architectures are configured to match in total
parameter count to within a small fraction. At the tiny scale, both
contain roughly 1.45\,M parameters (gap of 0.5\%). At the medium scale,
the xLSTM contains 4.55\,M parameters and the Transformer 4.22\,M (gap
of 7\%). The Transformer's MLP expansion factor is reduced from the
canonical~4 to~2 at the medium scale to keep the parameter budget
balanced. The complete architecture specification for each
configuration is summarised in Table~\ref{tab:archspec}.

\begin{table}[h]
\centering
\caption{Per-configuration architecture specification. \emph{Inner} for
  xLSTM denotes the (sLSTM, mLSTM) split inside the block stack; for the
  Transformer it denotes the MLP expansion factor relative to the
  embedding dimension.}
\label{tab:archspec}
\small
\begin{tabular}{lrrrrrr}
\toprule
Model & Vocab & Embed & Blocks & Heads & Context & Inner \\
\midrule
xLSTM tiny         & 4{,}000 & 128 & 4 (3$\times$mLSTM, 1$\times$sLSTM) & 4 & 256 & FFN $\times 1.3$ \\
xLSTM medium       & 4{,}000 & 256 & 6 (5$\times$mLSTM, 1$\times$sLSTM) & 4 & 256 & FFN $\times 1.3$ \\
Transformer tiny   & 4{,}000 & 128 & 2 & 4 & 256 & MLP $\times 4$ \\
Transformer medium & 4{,}000 & 256 & 4 & 8 & 256 & MLP $\times 2$ \\
\bottomrule
\end{tabular}
\end{table}

\paragraph{FLOPs equivalence.}
Parameter count is one budget; compute is another. Using the standard
Chinchilla approximation $\text{FLOPs}_{\text{fwd}} \approx 2N$ per token
(with $N$ the non-embedding parameter count) and adding the explicit
$\mathcal{O}(L^2 d)$ attention cost for the Transformer, we obtain
forward-pass costs of $1.90$\,M (xLSTM tiny), $2.08$\,M (Transformer
tiny), $7.06$\,M (xLSTM medium) and $7.31$\,M (Transformer medium) FLOPs
per token. The two architectures are matched in compute to within $9\%$
at the tiny scale and $4\%$ at the medium scale. The 7\% parameter gap
at the medium scale therefore corresponds to a tighter $\sim 4\%$ FLOPs
gap once the Transformer's quadratic-in-context attention term is
accounted for, which mitigates one obvious objection to a strictly
parameter-matched comparison.

\subsection{Training}

The training loop is shared across architectures and identical in every
hyperparameter except for the architecture itself. We use AdamW with
\(\beta_1{=}0.9\), \(\beta_2{=}0.95\), weight decay 0.1 (applied only
to tensors of rank $\ge$ 2, following the Chinchilla
convention~\cite{hoffmann2022}), peak learning rate \(3 \times 10^{-4}\),
linear warmup followed by cosine decay to 10\% of the peak, and gradient
clipping at norm 1.0. Sequences of length 256 are packed without padding,
using \texttt{</s>} as document separator. Batches contain 4 sequences;
no gradient accumulation is used at this scale. Training runs for 1{,}500
steps at the tiny scale and 2{,}000 steps at the medium scale. All
experiments use full \texttt{float32} precision on CPU. The complete
hyperparameter set is given in Table~\ref{tab:hyperparams}.

\begin{table}[h]
\centering
\caption{Training hyperparameters. All values are identical between
  xLSTM and Transformer at a given scale; the architecture identifier is
  the only differing field across paired runs.}
\label{tab:hyperparams}
\small
\begin{tabular}{lrr}
\toprule
Hyperparameter & Tiny & Medium \\
\midrule
Optimizer                 & \multicolumn{2}{c}{AdamW} \\
$(\beta_1, \beta_2)$      & \multicolumn{2}{c}{$(0.9,\, 0.95)$} \\
Weight decay (rank $\ge 2$) & \multicolumn{2}{c}{$0.1$} \\
Peak learning rate        & \multicolumn{2}{c}{$3\times 10^{-4}$} \\
Min LR (cosine floor)     & \multicolumn{2}{c}{$10\%$ of peak} \\
LR schedule               & \multicolumn{2}{c}{linear warmup, cosine decay} \\
Gradient clip (L2 norm)   & \multicolumn{2}{c}{$1.0$} \\
Sequence length           & \multicolumn{2}{c}{$256$} \\
Batch size                & \multicolumn{2}{c}{$4$} \\
Gradient accumulation     & \multicolumn{2}{c}{$1$} \\
Precision                 & \multicolumn{2}{c}{\texttt{float32}} \\
Random seed               & \multicolumn{2}{c}{$42$} \\
\midrule
Warmup steps              & $100$  & $200$ \\
Total steps               & $1{,}500$ & $2{,}000$ \\
Eval interval (steps)     & $100$  & $200$ \\
\bottomrule
\end{tabular}
\end{table}

\paragraph{Token budget.}
The training budgets correspond to $1.54$\,M tokens (tiny) and $2.05$\,M
tokens (medium). Expressed as the token-to-parameter ratio used in
Chinchilla-style analyses~\cite{hoffmann2022}, this places our runs at
ratios of $\approx 1.0$\,(tiny) and $\approx 0.45$\,(medium) tokens per
parameter, $20$ to $45\times$ \emph{below} the compute-optimal ratio of
$\sim 20{:}1$. Both runs are therefore deeply data-limited rather than
capacity-limited. We do not view this as an experimental flaw but as a
deliberate property of the regime studied: the corpus size is the
binding constraint imposed by the public-domain Portuguese literary
record. We expect this regime to favour architectures with weaker
inductive biases (i.e., fewer parameters spent on learnable gating
machinery), an observation we return to in the Discussion.

\subsection{Evaluation}

Evaluation is performed at fixed intervals during training on a
deterministic, file-based holdout dataset (Section~3.1). Tokenisation,
chunking order, and the number of evaluation chunks are fixed across
runs, so reported perplexity values are directly comparable. We report
the \emph{best} dev perplexity and the \emph{final} dev perplexity, the
former because it is the value used to select checkpoints in practice,
and the latter because it characterises end-of-training behaviour.

% =============================================================================
\section{Results}

\subsection{Held-out perplexity}\label{sec:held-out-ppl}

Table~\ref{tab:results} reports best and final dev perplexity along with
training throughput for the four configurations.
Table~\ref{tab:ppltrajectory} reports the held-out perplexity at fixed
training steps for the same four runs, so that the convergence behaviour
can be inspected at the same temporal grid across configurations.

\begin{table}[h]
\centering
\caption{Held-out perplexity on \emph{Memórias Póstumas de Brás Cubas}.
  Throughput is the mean training throughput observed during the run on
  CPU.}
\label{tab:results}
\small
\begin{tabular}{lrrrr}
\toprule
Model & Params & Best PPL & Final PPL & Throughput \\
\midrule
xLSTM tiny         & 1.46\,M & 461.9 & 514.7 & 2.8\,k tok/s \\
Transformer tiny   & 1.45\,M & \textbf{332.5} & 510.1 & 22.4\,k tok/s \\
xLSTM medium       & 4.55\,M & \textbf{229.9} & 229.9 & 1.8\,k tok/s \\
Transformer medium & 4.22\,M & 235.6 & 235.6 & 8.6\,k tok/s \\
\bottomrule
\end{tabular}
\end{table}

Three patterns are visible in Table~\ref{tab:ppltrajectory}: (i)~the
convergence is monotone with intermittent transient spikes (e.g.\ xLSTM
tiny at step~300, both tiny runs at step~900 --- the latter coincides
with the cosine decay beginning); (ii)~at every checkpoint past
step~600, the medium models dominate their tiny counterparts; (iii)~the
medium-scale crossover between architectures emerges only past
step~1200, with the xLSTM and Transformer crossing one another twice
(steps~1400 and~1800) before settling at step~2000.

\begin{table}[h]
\centering
\caption{Held-out dev perplexity at fixed training steps. Best-per-row
  values are in bold. ``--'' marks steps past the run's training budget.}
\label{tab:ppltrajectory}
\small
\setlength{\tabcolsep}{4pt}
\begin{tabular}{lrrrr}
\toprule
Step & xLSTM tiny & Transformer tiny & xLSTM medium & Transformer medium \\
\midrule
200   & 947.7 & 1184.7 & 1200.3 & 1386.8 \\
400   & 868.6 & 678.5  & 730.7  & 690.3  \\
600   & 607.2 & 435.4  & 414.0  & 400.8  \\
800   & 553.1 & 399.0  & 392.0  & 405.9  \\
1000  & 566.7 & 388.1  & 355.6  & 380.8  \\
1200  & 480.1 & 342.7  & 278.3  & 286.0  \\
1400  & \textbf{461.9} & \textbf{332.5}  & 260.2  & 279.9 \\
1500  & 514.7 & 510.1  & --     & --     \\
1600  & --    & --     & 280.3  & 293.6  \\
1800  & --    & --     & 235.7  & 243.9  \\
2000  & --    & --     & \textbf{229.9}  & \textbf{235.6} \\
\bottomrule
\end{tabular}
\end{table}

First, both architectures benefit substantially from the move from
tiny to medium scale: best held-out perplexity is reduced by approximately
50\% in both cases. This indicates that the tiny scale is undertrained
relative to the architectures' capacity, and that the held-out perplexity
is genuinely sensitive to model capacity in this regime.

Second, the relative ranking between architectures inverts with scale.
At the tiny scale, the Transformer attains a substantially lower best
perplexity than the xLSTM (332.5 vs. 461.9, a 28\% gap). At the medium
scale, the ordering inverts and the xLSTM marginally outperforms the
Transformer (229.9 vs. 235.6, a 2.5\% gap). The medium-scale gap is
small enough that, in the absence of multiple training seeds, no firm
claim of superiority can be made; we discuss this further in
Section~\ref{sec:limitations}.

A subsidiary observation is that the xLSTM medium run reached its best
held-out perplexity at the final evaluation, indicating that the run was
\emph{undertrained} at the budget chosen. The Transformer medium also
ended with best equal to final, suggesting both runs would benefit from
additional steps; a confirmatory run with the longer schedule is left
to future work.

\subsection{Sample efficiency}\label{sec:sample-efficiency}

The trajectory data also lets us re-frame the comparison in terms of
\emph{tokens to reach a target perplexity} rather than perplexity at a
fixed step count. Table~\ref{tab:tokens-to-ppl} reports, for each run,
the smallest number of training tokens required to reach a sequence of
perplexity thresholds. This view is more directly informative about
sample efficiency than endpoint perplexity, which conflates capacity
with optimisation budget.

\begin{table}[h]
\centering
\caption{Training tokens required to first reach a given held-out
  perplexity threshold (lower is more sample-efficient). ``--'' marks
  thresholds the run never reaches within its budget.}
\label{tab:tokens-to-ppl}
\small
\setlength{\tabcolsep}{6pt}
\begin{tabular}{lrrrr}
\toprule
PPL target & xLSTM tiny & Transformer tiny & xLSTM medium & Transformer medium \\
\midrule
$\le 1{,}000$ & $0.20$\,M & $0.41$\,M & $0.41$\,M & $0.41$\,M \\
$\le 600$     & $0.72$\,M & $0.51$\,M & $0.61$\,M & $0.61$\,M \\
$\le 400$     & --        & $0.72$\,M & $0.82$\,M & $1.02$\,M \\
$\le 300$     & --        & --        & $1.23$\,M & $1.23$\,M \\
$\le 250$     & --        & --        & $1.84$\,M & $1.84$\,M \\
\bottomrule
\end{tabular}
\end{table}

Two observations: (i)~at the tiny scale the Transformer crosses every
threshold first that both architectures eventually reach, consistent
with its lower endpoint perplexity; (ii)~at the medium scale the two
architectures reach the same thresholds at the same step counts within
the resolution of our eval grid (every 200 steps), and the xLSTM's
final-step advantage is driven by the very last evaluations rather
than by faster crossing of any earlier threshold. Sample-efficiency
parity at medium scale, combined with end-of-training divergence, is
also consistent with the \emph{undertrained} reading: if both models
were trained past the cosine-decay tail, the gap could plausibly grow,
shrink, or invert.

\subsection{Quantitative analysis of generated text}\label{sec:gentext}

To complement the qualitative samples, we compute two surface-level
statistics on \emph{generated} (not held-out) text from the medium-scale
models: type-token ratio (TTR) and verbatim 4-gram repetition rate. We
sample 5 seeds $\times$ 4 prompts $\times$ 120 new tokens per
configuration (matched decoding settings: temperature 0.8, top-$k$ 40)
and concatenate. As a reference we report the same statistics on a
length-matched 1{,}500-word slice of the held-out corpus
(\emph{Brás Cubas}). Results are in Table~\ref{tab:gentext}.

\begin{table}[h]
\centering
\caption{Surface statistics on $\sim$1.2--1.4\,k generated whitespace
  tokens. Type-token ratio (TTR) measures lexical diversity; 4-gram
  repetition rate is the fraction of 4-grams that have already appeared
  earlier in the same generated stream. Reference values come from a
  1{,}500-word slice of the held-out novel.}
\label{tab:gentext}
\small
\begin{tabular}{lrrr}
\toprule
Source & Tokens & TTR & 4-gram rep. \\
\midrule
xLSTM medium (generated)       & $1{,}416$ & $0.277$ & $0.188$ \\
Transformer medium (generated) & $1{,}192$ & $0.302$ & $0.080$ \\
\midrule
Reference (\emph{Brás Cubas}, 1.5\,k words) & $1{,}500$ & $0.564$ & $0.002$ \\
\bottomrule
\end{tabular}
\end{table}

The two architectures differ in a direction not visible from
perplexity: the Transformer produces measurably more diverse and less
repetitive output than the xLSTM ($+9\%$ relative TTR, less than
half the 4-gram repetition rate), despite a marginally \emph{worse}
held-out perplexity. Both models remain far from the human-text
references on both metrics, as expected at this scale. We flag the
divergence between perplexity ranking and generation-time repetition as
a candidate signal for future work: surface repetition is a known
failure mode of small autoregressive LMs, and it is not obvious that
the architecture which scores better in held-out PPL also generates
better samples on these surface measures.

\subsection{Downstream task: author classification probe}\label{sec:probe}

Held-out perplexity measures how well the model predicts the next token,
but it conflates lexical and stylistic signal in a way that may not
match what a downstream user would care about. We complement
Section~\ref{sec:held-out-ppl} with a standard linear-probe protocol:
freeze the trained LM, extract mean-pooled hidden states from
fixed-length passages, and train a single linear classifier to predict
the author of each passage. The probe eval gives an architecture-level
comparison metric that does not depend on perplexity calibration.

\paragraph{Protocol.}
For every novel in the training corpus we cut non-overlapping passages
of $128$ tokens (using the same byte-level BPE tokenizer the LMs were
trained with), capping at $400$ passages per author to control class
imbalance. This yields $9$ authors $\times\, 400 = 3{,}600$ passages.
Each passage is fed through the frozen LM; we mean-pool the
hidden-state output over the time axis (i.e., the output of
\texttt{xlstm\_block\_stack} for xLSTM, or the post-final-norm hidden
state for the Transformer) to obtain one $D$-dimensional feature vector.
Features are standardised per dimension using train-split statistics. A
single \texttt{nn.Linear} layer is trained for $200$ epochs with AdamW
($\eta = 10^{-2}$, weight decay $10^{-3}$), cross-entropy loss, on a
stratified $80\%/20\%$ split. We report accuracy and macro-averaged F1
on the held-out $20\%$.

\paragraph{Out-of-pretraining authors.}
The pretraining corpus that produced the medium-scale checkpoints used
six novels from three authors (Machado de Assis, José de Alencar, Eça
de Queirós). The probe is run over the expanded $9$-author corpus that
the models have \emph{not} seen during pretraining for six of the nine
authors. The probe therefore tests whether the LMs' representations
support author identification both for in-pretraining and
out-of-pretraining authors.

\paragraph{Results.}
Table~\ref{tab:probe} reports overall and per-author metrics for both
medium-scale checkpoints. Both architectures clear the
majority-class baseline ($11.1\%$, since the per-author cap balances
the dataset) by a wide margin, indicating that the LMs encode
substantial author-discriminative signal even though they were not
trained for classification. The xLSTM outperforms the Transformer by
$6.7$ accuracy points and $6.9$ macro-F1 points; this ranking is in the
same direction as the held-out-perplexity ranking at the medium scale
(xLSTM wins both), which is some evidence that the perplexity advantage
is not an artefact of the metric. Both architectures are strongest on
authors that contributed multiple novels to pretraining (José de
Alencar, Machado de Assis) and on authors with the most stylistically
distinctive prose in the set (Bernardo Guimarães, Eça de Queirós), and
weakest on the most prolific author of the period in general (Camilo
Castelo Branco) and on Júlio Dinis, both unseen during pretraining and
both Portuguese realists with overlapping literary register.

\begin{table}[h]
\centering
\caption{Author-classification linear probe on $128$-token passages
  (9 authors, $400$ passages each, $80\%/20\%$ stratified split).
  Bold marks the higher value per row.}
\label{tab:probe}
\small
\setlength{\tabcolsep}{6pt}
\begin{tabular}{lrr}
\toprule
Metric & xLSTM medium & Transformer medium \\
\midrule
Test accuracy             & \textbf{66.4\%} & 59.7\% \\
Test macro-F1             & \textbf{66.4\%} & 59.5\% \\
Train accuracy (sanity)   & $83.3\%$        & $73.3\%$ \\
Majority-class baseline   & \multicolumn{2}{c}{$11.1\%$} \\
\midrule
\multicolumn{3}{l}{\emph{Per-author F1 (test set):}} \\
Almeida Garrett           & \textbf{60.7\%} & 54.9\% \\
Aluísio Azevedo           & \textbf{54.1\%} & 53.2\% \\
Bernardo Guimarães        & \textbf{70.5\%} & 64.6\% \\
Camilo Castelo Branco     & \textbf{54.4\%} & 51.5\% \\
Eça de Queirós            & \textbf{70.7\%} & 56.9\% \\
José de Alencar           & \textbf{86.5\%} & 81.3\% \\
Júlio Dinis               & \textbf{52.9\%} & 42.1\% \\
Machado de Assis          & \textbf{83.4\%} & 74.4\% \\
Raul Pompéia              & \textbf{64.2\%} & 56.8\% \\
\bottomrule
\end{tabular}
\end{table}

\paragraph{Caveats.}
The probe is trained on top of frozen features; we do not fine-tune the
LM. Per-author macro-F1 numbers should be read in the light of the
small per-class test set ($80$ passages) and the absence of multi-seed
estimates for the probe itself, which we discuss in
Section~\ref{sec:limitations}.

\subsection{Downstream task: sentiment probe (IMDB pt-BR)}\label{sec:sentiment}

The author probe of Section~\ref{sec:probe} measures how well the LM
encodes the \emph{stylistic} signal that distinguishes 19th-century
authors --- a signal close to the LM's pretraining distribution.
To complement it with a task that is genuinely \emph{out-of-distribution}
for our pretraining corpus, we run the same linear-probe protocol on a
binary sentiment dataset: \texttt{celsowm/imdb-reviews-pt-br}, the
$\sim$49\,k IMDB review collection translated to Portuguese.\footnote{%
\url{https://huggingface.co/datasets/celsowm/imdb-reviews-pt-br}.
The original English IMDB reviews of Maas et
al.~\cite{maas2011imdb}, machine-translated to Brazilian Portuguese.}
The pretraining corpus contains no movie reviews and no contemporary
Portuguese: a strong probe score under this shift is evidence that the
LM has learned representations that transfer beyond the
narrow-domain training data.

\paragraph{Protocol.}
We sample $2{,}000$ reviews per class ($4{,}000$ in total) from the
dataset, tokenise each one with the byte-level BPE that was used for
pretraining (Section~3.2), truncate or right-pad to $256$ tokens, and
extract one feature vector per review by mean-pooling the frozen LM's
hidden states over the time axis. The features are standardised with
train-split statistics, stratified-split $80\%/20\%$, and used to train
a single \texttt{nn.Linear(D, 2)} layer for $200$ epochs with AdamW
($\eta = 10^{-2}$, weight decay $10^{-3}$). The protocol is identical
to the author probe except for (i)~the dataset, (ii)~the number of
classes, and (iii)~the per-review truncation in place of the per-book
sliding window.

\paragraph{Results.}
Table~\ref{tab:sentiment} reports the metrics. Both architectures
clear the $50\%$ majority-class baseline; the xLSTM outperforms the
Transformer by $7.9$ accuracy points and $7.9$ macro-F1 points on this
out-of-domain task. The ranking direction is the same as the author
probe (Table~\ref{tab:probe}) and as held-out perplexity at the medium
scale (Table~\ref{tab:results}). We caveat this finding in
Section~\ref{sec:olist}, where a multi-seed re-run of the same protocol
on a native-Brazilian-Portuguese review dataset shows the two
architectures within one standard deviation of each other.

\begin{table}[h]
\centering
\caption{Binary sentiment linear probe on
  \texttt{celsowm/imdb-reviews-pt-br} ($2{,}000$ reviews per class,
  $256$-token features, $80\%/20\%$ stratified split). Bold marks the
  higher value per row.}
\label{tab:sentiment}
\small
\setlength{\tabcolsep}{6pt}
\begin{tabular}{lrr}
\toprule
Metric & xLSTM medium & Transformer medium \\
\midrule
Test accuracy            & \textbf{67.4\%} & 59.5\% \\
Test macro-F1            & \textbf{67.4\%} & 59.5\% \\
Train accuracy (sanity)  & $70.2\%$        & $66.2\%$ \\
Per-class F1 (negative)  & \textbf{68.1\%} & 59.9\% \\
Per-class F1 (positive)  & \textbf{66.7\%} & 59.1\% \\
Majority-class baseline  & \multicolumn{2}{c}{$50.0\%$} \\
\bottomrule
\end{tabular}
\end{table}

\paragraph{Caveats.}
Both train and test accuracies are far below the $\sim$90\% routinely
achieved on English IMDB by full transformer encoders; this is
expected. The pretraining corpus consists of $\sim$1.9\,M tokens of
literary 19th-century Portuguese, vocabulary $4{,}000$, and the LMs
have $4$--$5$\,M parameters --- two to four orders of magnitude below
the size and three corpora removed from the domain that produce
state-of-the-art sentiment numbers. The probe is therefore best read
as an architecture-comparison signal under shift, not as an absolute
sentiment-classification result.

\subsection{Downstream task: review-score probe (Olist e-commerce, native PT-BR)}\label{sec:olist}

The IMDB-pt probe of Section~\ref{sec:sentiment} is a useful
out-of-distribution check, but the data is machine-translated from
English; any artefact of the MT system enters the probe along with the
sentiment signal. To remove that confound and to add error bars to the
single-seed numbers reported in Table~\ref{tab:sentiment}, we re-run
the same probe protocol on the Olist Brazilian E-Commerce Public
dataset, using the free-text customer reviews
(\texttt{review\_comment\_message}) as inputs and the $1$--$5$ star
rating (\texttt{review\_score}) as the label.\footnote{%
\url{https://www.kaggle.com/datasets/olistbr/brazilian-ecommerce}.
Native Brazilian-Portuguese reviews from a real Brazilian e-commerce
platform; no translation is involved.} The dataset is genuinely
in-language but, like IMDB, fully out-of-domain for our literary
pretraining corpus, and the messages are short and informal --- a
register the LM has never seen.

\paragraph{Protocol.}
Of the $99{,}224$ rows in the dataset, $40{,}950$ have a non-empty
review message; we keep only those. To match the binary IMDB protocol
we drop reviews with \texttt{review\_score}$=3$ (the ambiguous middle)
and label $1$--$2$ as \emph{negative} and $4$--$5$ as \emph{positive}.
We sample $2{,}000$ reviews per class ($4{,}000$ in total),
tokenise with the same byte-level BPE used in pretraining, and
truncate or right-pad to $128$ tokens. The shorter window matches
the data: review messages have a median length of $53$ characters
(maximum $208$), so $128$ tokens suffice. Feature extraction,
standardisation, $80\%/20\%$ stratified split, and the
\texttt{nn.Linear(D, 2)} probe are identical to
Section~\ref{sec:sentiment}. The single methodological change is
\emph{three independent seeds} ($42$, $1337$, $2024$); each seed
re-samples the data, re-splits, and re-initialises the probe, and we
report mean $\pm$ standard deviation across seeds for every metric.

\paragraph{Results.}
Table~\ref{tab:olist} reports the metrics. Both architectures clear
the $50\%$ majority-class baseline by $\sim$33 points and end up
within one standard deviation of each other on test accuracy and
macro-F1. The point estimate places the Transformer $0.83$ accuracy
points above the xLSTM, with $1\sigma$ intervals
$[81.9, 83.9]$ vs.\ $[82.9, 84.7]$ that overlap throughout. The
ordering on this dataset is therefore opposite to the IMDB-pt
result of Section~\ref{sec:sentiment}, although in both cases the
xLSTM-vs-Transformer gap on the Olist task is well inside the
seed-to-seed noise we now measure. We read the combination as
follows: the IMDB-pt advantage of the xLSTM (a $7.9$-point gap from a
single seed) is real for that translated dataset, but the multi-seed
Olist result shows that on a native-Portuguese, in-language
out-of-domain probe the two architectures are statistically
indistinguishable at this scale. The ``three scalars agree'' claim
of Section~\ref{sec:sentiment} should therefore be read as
\emph{three out of four}: held-out perplexity, the author probe, and
the IMDB-pt sentiment probe rank xLSTM ahead; the multi-seed Olist
review probe records a tie.

\begin{table}[h]
\centering
\caption{Review-score linear probe on Olist e-commerce reviews
  ($2{,}000$ reviews per class, $128$-token features, binary
  $1$--$2$ vs.\ $4$--$5$, $80\%/20\%$ stratified split). All numbers
  are mean $\pm$ standard deviation across three seeds
  ($42$, $1337$, $2024$); each seed is a full end-to-end re-run.}
\label{tab:olist}
\small
\setlength{\tabcolsep}{6pt}
\begin{tabular}{lrr}
\toprule
Metric & xLSTM medium & Transformer medium \\
\midrule
Test accuracy            & $82.92 \pm 1.01\%$ & $\mathbf{83.75 \pm 0.90\%}$ \\
Test macro-F1            & $82.91 \pm 1.01\%$ & $\mathbf{83.75 \pm 0.90\%}$ \\
Train accuracy (sanity)  & $83.83 \pm 0.63\%$ & $85.35 \pm 0.24\%$ \\
Per-class F1 (negative)  & $82.51 \pm 0.91\%$ & $\mathbf{83.65 \pm 0.99\%}$ \\
Per-class F1 (positive)  & $83.30 \pm 1.11\%$ & $\mathbf{83.85 \pm 0.82\%}$ \\
Feature-extraction time  & $174.2 \pm 4.9$\,s & $\mathbf{25.8 \pm 0.7}$\,s \\
Throughput (passages/s)  & $23.0$            & $\mathbf{155.0}$ \\
Majority-class baseline  & \multicolumn{2}{c}{$50.0\%$} \\
\bottomrule
\end{tabular}
\end{table}

The accuracy columns are within one standard deviation, but the
feature-extraction columns differ by a factor of $\sim$$6.7\times$ on
this CPU benchmark. We discuss this gap in
Section~\ref{sec:throughput}: it is a property of our
\texttt{vanilla}-backend xLSTM running on CPU at short sequence
lengths, not a fundamental architectural cost.

\paragraph{Caveats.}
The same caveats as for the IMDB-pt probe apply: the LMs are small
($\sim$$4$--$5$\,M parameters), the pretraining corpus is literary,
and review messages are short and colloquial. The probe should be
read as an architecture-comparison signal, not as a competitive
sentiment classifier on Olist. One additional caveat is specific to
this dataset: the $1$--$5$ star distribution is heavily skewed
towards $5$ ($\sim$$50\%$ of reviews with text), and our binary cut
discards the $\sim$$8.7\%$ of $3$-star reviews. A multi-class
$1$--$5$ extension is supported by the script (\texttt{--multiclass}
flag) and would let one ask whether the representations capture
review intensity beyond polarity, but we leave that to future work.

\subsection{Throughput and the implementation caveat}\label{sec:throughput}

Table~\ref{tab:results} also reports the mean training throughput in
tokens per second. The Transformer trains roughly $5$--$8\times$ faster
than the xLSTM in our CPU setting; on GPU, the gap widens to
$\sim$$45$--$80\times$ (see Section~\ref{sec:gpu-scaling} and
Table~\ref{tab:gpu-results}). We emphasise that this is \emph{not} a
fundamental architectural cost. The reported numbers reflect the use of
the \texttt{vanilla} (pure PyTorch) xLSTM backend; the dedicated CUDA
kernels of the official xLSTM implementation, which parallelise the
recurrence via a custom scan operation, would close most of this gap on
GPU. Conversely, the Transformer benefits substantially in CPU from the
fused \texttt{F.scaled\_dot\_product\_attention} kernel, which exploits
SIMD/AVX in a way that the loop-based recurrent path does not; on GPU
the same kernel dispatches to FlashAttention~\cite{dao2022flashattention},
amplifying the effect.

This implementation-versus-architecture distinction is a methodological
point that the pipeline makes explicit by logging throughput in the same
\texttt{log.jsonl} as the perplexity numbers, and that we believe is
under-reported in arch-comparison papers in general. We return to it
quantitatively in Section~\ref{sec:gpu-scaling}, where the GPU
throughput numbers make the implementation-versus-architecture gap
$\sim$$10\times$ larger than it appears on CPU.

\subsection{GPU-scale architecture comparison}\label{sec:gpu-scaling}

To extend the comparison beyond the CPU regime studied in
Sections~\ref{sec:held-out-ppl}--\ref{sec:olist}, we ran \emph{both}
architectures on a single consumer-grade GPU (NVIDIA RTX~5070~Ti,
Compute Capability $12.0$) at three parameter scales: $\sim$$2.5$\,M,
$\sim$$6.3$--$6.6$\,M, and $\sim$$10$\,M. The GPU runs use a larger
vocabulary ($8{,}000$ instead of $4{,}000$), wider context at the
$10$\,M scale ($512$ instead of $256$), an $8\times$ larger batch
($32$ vs.\ $4$), and proportionally larger token budgets ($1{,}500$,
$4{,}000$, $8{,}000$ steps). The tokenizer differs from the CPU
experiments, so the GPU perplexities are not directly comparable to
Table~\ref{tab:results}; the runs form a self-contained GPU sweep.

\paragraph{Held-out perplexity.}
Table~\ref{tab:gpu-results} reports best and final held-out
perplexity together with mean training throughput.
Figure~\ref{fig:gpu-scaling} shows the training trajectories. Three
observations follow. First, the Transformer attains the lowest best
perplexity at the $10$\,M scale ($269.3$ vs.\ xLSTM $287.3$, a $6\%$
relative advantage). Second, the cross-architecture ordering is
\emph{non-monotone in scale}: the Transformer wins at tiny by a
$23\%$ margin, the xLSTM marginally surpasses the Transformer at
medium by $9\%$, and the Transformer regains the lead at $10$\,M.
Third, the xLSTM $10$\,M run reaches its best perplexity at step
$1{,}600$ of $8{,}000$ ($20\%$ into the schedule) and drifts
upwards by $53\%$ thereafter, while the Transformer $10$\,M run
keeps improving until step $4{,}400$ and finishes with a far milder
final-vs-best ratio ($1.20$ vs.\ $1.53$). The xLSTM at $10$\,M is
therefore the only configuration in our sweep that visibly overfits
the corpus. Figure~\ref{fig:best-ppl-vs-params} plots the best dev
perplexity as a function of parameter count, making the two
crossovers and the relative compactness of the $10$\,M gap easier
to read off than from Table~\ref{tab:gpu-results} alone.

\begin{table}[h]
\centering
\caption{GPU-scale architecture comparison. Tokens-seen reflects the
  training budget (steps $\times$ batch $\times$ sequence length).
  Throughput is the mean tokens/s observed during training, with the
  first measurement excluded as warmup. Best-of-pair values per row
  are in bold.}
\label{tab:gpu-results}
\small
\setlength{\tabcolsep}{4pt}
\begin{tabular}{lrrrrrr}
\toprule
Model & Params & Steps & Tokens & Best PPL & Final PPL & Throughput \\
\midrule
xLSTM tiny          & $2.5$\,M  & $1{,}500$ & $12.3$\,M  & $699.5$ & $733.9$ & $10.3$\,k tok/s \\
Transformer tiny    & $2.5$\,M  & $1{,}500$ & $12.3$\,M  & $\mathbf{539.3}$ & $814.7$ & $\mathbf{702.1}$\,k tok/s \\
xLSTM medium        & $6.6$\,M  & $4{,}000$ & $32.8$\,M  & $\mathbf{377.7}$ & $\mathbf{405.4}$ & $10.1$\,k tok/s \\
Transformer medium  & $6.3$\,M  & $4{,}000$ & $32.8$\,M  & $413.0$ & $418.4$ & $\mathbf{450.5}$\,k tok/s \\
xLSTM $10$\,M       & $10.1$\,M & $8{,}000$ & $131.1$\,M & $287.3$ & $440.3$ & $5.9$\,k tok/s \\
Transformer $10$\,M & $10.2$\,M & $8{,}000$ & $131.1$\,M & $\mathbf{269.3}$ & $\mathbf{322.4}$ & $\mathbf{471.2}$\,k tok/s \\
\bottomrule
\end{tabular}
\end{table}

\begin{figure}[h]
\centering
\includegraphics[width=\linewidth]{gpu_arch_curve.png}
\caption{Training loss (left) and held-out perplexity (right, log
  scale) for the six GPU-scale runs (xLSTM and Transformer at
  $2.5$\,M, $\sim$$6.5$\,M, and $\sim$$10$\,M parameters). The
  $10$\,M xLSTM trajectory (green) is the only one in which the
  held-out perplexity diverges \emph{upward} from the training loss
  past step $\sim$$2{,}000$ --- the textbook overfitting signature.}
\label{fig:gpu-scaling}
\end{figure}

\begin{figure}[h]
\centering
\includegraphics[width=0.78\linewidth]{best_ppl_vs_params.png}
\caption{Best held-out perplexity as a function of parameter count for
  the GPU sweep. The two architectures cross twice across the swept
  range: the Transformer leads at $2.5$\,M, the xLSTM leads at
  $\sim$$6.5$\,M, the Transformer regains the lead at $10$\,M. The
  shaded band marks the medium-scale crossover.}
\label{fig:best-ppl-vs-params}
\end{figure}

\paragraph{Throughput on GPU: an $\sim$$80\times$ gap.}
The Throughput column of Table~\ref{tab:gpu-results} exposes a
\emph{$\sim$$45$--$80\times$} speed gap between architectures on
GPU. At the $10$\,M scale the Transformer processes $471.2$\,k
tokens/s while the xLSTM processes $5.9$\,k tokens/s, a ratio of
$80\times$; at the medium and tiny scales the ratios are
$45\times$ and $68\times$ respectively. This dwarfs the
$5$--$8\times$ CPU gap of Section~\ref{sec:throughput}. The
asymmetry is implementation-driven, not architectural: PyTorch
dispatches \texttt{F.scaled\_dot\_product\_attention} to
FlashAttention~\cite{dao2022flashattention} on supported
GPUs --- a memory-IO-aware, kernel-fused attention path against
which our \texttt{vanilla} pure-PyTorch xLSTM backend has no
equivalent. The dedicated CUDA kernels of the official xLSTM
release~\cite{beck2024} would close most of this gap but require a
higher Compute Capability than the CUDA build available to us. The
methodological point of Section~\ref{sec:throughput} is therefore
\emph{amplified}, not weakened, on GPU: at the scales we study, a
$6\%$ perplexity advantage in favour of the Transformer comes
alongside a $\sim$$80\times$ wall-clock advantage for the
Transformer, and the two numbers must be read together.
Figure~\ref{fig:throughput} visualises the gap on a log scale.

\begin{figure}[h]
\centering
\includegraphics[width=0.80\linewidth]{throughput_bars.png}
\caption{GPU training throughput (mean tokens/s, log scale) per
  architecture and scale on the NVIDIA RTX~5070~Ti. Ratios above the
  Transformer bars give the per-scale speedup
  (Transformer / xLSTM\,vanilla). The xLSTM throughput drops from
  $10$\,k to $5.9$\,k tok/s between medium and $10$\,M because the
  $10$\,M configuration doubles the context length and adds a second
  sLSTM block.}
\label{fig:throughput}
\end{figure}

\paragraph{Downstream: author probe on GPU checkpoints.}
We rerun the linear-probe protocol of Section~\ref{sec:probe} on
the medium and $10$\,M GPU checkpoints. The dataset, protocol
($9$ authors, $400$ passages each, $128$-token windows, $80/20$
stratified split, mean-pooled features, \texttt{nn.Linear}
probe), and majority baseline ($11.1\%$) are unchanged.
Table~\ref{tab:probe-gpu} reports accuracy and macro-F1.

\begin{table}[h]
\centering
\caption{Author-classification linear probe on the GPU-scale
  checkpoints (same protocol as Table~\ref{tab:probe}; majority
  baseline $11.1\%$). Bold marks the higher value per scale.}
\label{tab:probe-gpu}
\small
\setlength{\tabcolsep}{6pt}
\begin{tabular}{lrr}
\toprule
Model & Test acc. & Test macro-F1 \\
\midrule
xLSTM medium        & $\mathbf{66.11\%}$ & $\mathbf{66.07\%}$ \\
Transformer medium  & $64.17\%$          & $63.95\%$          \\
\midrule
xLSTM $10$\,M       & $\mathbf{68.06\%}$ & $\mathbf{67.77\%}$ \\
Transformer $10$\,M & $65.83\%$          & $65.71\%$          \\
\bottomrule
\end{tabular}
\end{table}

The xLSTM leads at both medium and $10$\,M scales (by $1.94$ and
$2.23$ accuracy points respectively). The direction matches the
CPU result of Table~\ref{tab:probe} but the magnitude is smaller,
consistent with the larger GPU training budget reducing the
probe-level architectural gap. Moving from medium to $10$\,M
improves both architectures, indicating the extra capacity captures
additional author-discriminative signal. Note that at $10$\,M the
perplexity ranking (Transformer better) and the probe ranking
(xLSTM better) \emph{disagree} --- the same generation-vs-perplexity
divergence already observed on CPU
(Section~\ref{sec:gentext}).

\paragraph{Downstream: IMDB-PT sentiment via in-pipeline MarianMT translation.}
Section~\ref{sec:sentiment} used the pre-translated
\texttt{celsowm/imdb-reviews-pt-br}. To remove dependence on a
third-party translation and to obtain an end-to-end reproducible
sentiment evaluation, we add an in-pipeline translation step: we
download the original English IMDB
dataset~\cite{maas2011imdb} and translate a balanced subset
($1{,}000$ positive $+$ $1{,}000$ negative) from English to
Brazilian Portuguese using
\texttt{Helsinki-NLP/opus-mt-tc-big-en-pt}~\cite{tiedemann2020opus},
the OPUS-MT public en$\rightarrow$pt MarianMT~\cite{junczys2018marian}
checkpoint. OPUS-MT is the encoder--decoder Transformer family that
was deployed in production neural translation immediately after
GNMT and remains a strong public en$\rightarrow$pt baseline. The
translation is sentence-by-sentence with greedy decoding; the
resulting JSONL is then consumed by exactly the same probe protocol
as Section~\ref{sec:sentiment} ($256$-token features,
$80\%/20\%$ split, single \texttt{nn.Linear(D,2)}).
Table~\ref{tab:sentiment-gpu} reports the metrics on the GPU
$10$\,M checkpoints.

\begin{table}[h]
\centering
\caption{Binary sentiment linear probe on MarianMT-translated
  IMDB-PT ($2{,}000$ reviews balanced $50/50$, $256$-token features,
  $80\%/20\%$ stratified split). Bold marks the higher value per
  row. Majority-class baseline is $50\%$.}
\label{tab:sentiment-gpu}
\small
\setlength{\tabcolsep}{6pt}
\begin{tabular}{lrr}
\toprule
Metric & xLSTM $10$\,M & Transformer $10$\,M \\
\midrule
Test accuracy             & $\mathbf{64.00\%}$ & $63.25\%$ \\
Test macro-F1             & $\mathbf{64.00\%}$ & $63.24\%$ \\
Train accuracy (sanity)   & $75.06\%$          & $74.81\%$ \\
Per-class F1 (negative)   & $\mathbf{63.82\%}$ & $62.60\%$ \\
Per-class F1 (positive)   & $\mathbf{64.18\%}$ & $63.88\%$ \\
\bottomrule
\end{tabular}
\end{table}

The xLSTM leads by $0.75$ accuracy points and $0.76$ macro-F1
points. The gap is well inside the seed-to-seed band measured for
the Olist probe of Section~\ref{sec:olist} (where the two
architectures landed within $1\sigma$ of each other on a comparable
binary-sentiment task), so we read the result as \emph{statistical
parity} on out-of-domain binary sentiment at $10$\,M, with a thin
xLSTM lead. Combined with Table~\ref{tab:probe-gpu} and
Table~\ref{tab:gpu-results}, the $10$\,M picture is:
the Transformer reaches lower next-token perplexity ($6\%$ gap), but
the xLSTM's representations match or marginally surpass the
Transformer's on every downstream probe we measured, while training
$\sim$$80\times$ slower in wall-clock terms. Figure~\ref{fig:downstream}
summarises the downstream comparison visually. The in-pipeline
MarianMT translation step is, by itself, a useful methodological
addition: it lets us probe a domain (movie reviews) and a register
(contemporary, informal, conversational) that are absent from any
reasonably-sized Portuguese public-domain pretraining corpus,
without taking the translation system as a black-box dependency.

\begin{figure}[h]
\centering
\includegraphics[width=0.82\linewidth]{downstream_bars.png}
\caption{Linear-probe test accuracy of the GPU checkpoints on three
  tasks: $9$-author classification at medium and $10$\,M scales, and
  binary IMDB-PT sentiment classification at the $10$\,M scale. Dashed
  horizontal lines mark the majority-class baselines ($11.1\%$ for the
  $9$-author task, $50\%$ for binary sentiment). The xLSTM leads on
  all three tasks at the GPU scales considered, though the
  IMDB-PT gap is well inside the seed-to-seed band measured for the
  Olist probe of Section~\ref{sec:olist}.}
\label{fig:downstream}
\end{figure}

\paragraph{Why the $10$\,M xLSTM overfits early.}
At the corpus size used here ($\sim$$1.9$\,M unique tokens),
increasing xLSTM capacity past $\sim$$6$\,M parameters yields
diminishing held-out returns and eventually overfits within
$\sim$$25$\,M training tokens (Figure~\ref{fig:gpu-scaling}). The
Transformer $10$\,M run --- same parameter count, same data ---
overfits much less aggressively over the same budget. We interpret
this as consistent with the well-known empirical observation that
architectures rich in inductive bias (gates, matrix memory,
exponential gating with stabiliser state, GroupNorm head-wise) are
more prone to memorisation than attention-only architectures at
fixed corpus size, particularly when the token-to-parameter ratio is
deep below
Chinchilla-optimal~\cite{hoffmann2022}. A practical implication for
the pipeline user is that the $10$\,M xLSTM configuration would
benefit from either (i)~a smaller capacity given the corpus size,
(ii)~stronger regularisation (dropout, larger weight decay), or
(iii)~a substantially larger pretraining corpus; the CPU experiments
of Section~\ref{sec:held-out-ppl} live in regime (iii)-from-below
(under-trained) and the GPU $10$\,M xLSTM lives in regime
(iii)-from-above (over-capacity for the corpus).

\subsection{Qualitative inspection}

We provide samples from each trained model in Appendix~\ref{app:samples},
including a multi-prompt sweep over the medium-scale models that probes
behaviour beyond a single conditioning context.
Briefly, the medium-scale models produce text that is recognisably
Portuguese and reproduces the orthographic conventions of 19th-century
Brazilian literature (\emph{pae} for \emph{pai}, contractions like
\emph{n'um}, dialogue with em-dashes). The tiny-scale models produce
fragments with isolated correct words but limited syntactic coherence.
Sample-level impressions are subjective; the surface statistics in
Section~\ref{sec:gentext} (Table~\ref{tab:gentext}) provide the
quantitative companion to this observation.

% =============================================================================
\section{Discussion}

The pipeline succeeds in producing comparable runs in the sense that
matters for arch comparison: every component except the architecture is
shared, every hyperparameter is logged, and the held-out evaluation is
deterministic. We believe this is a useful artefact in itself,
independently of the specific result reported. In particular, the
pipeline can be extended to (i)~additional architectures by adding a
new branch to the model factory; (ii)~larger scales by adjusting the
configuration files (a 50\,M-parameter configuration is included in the
repository as a target for future GPU experiments); and (iii)~other
corpora by replacing the book-list module.

The crossover observation in our case study is consistent with the
hypothesis advanced by Beck et al.~\cite{beck2024} that xLSTM advantages
emerge with scale. We stress, however, that the gap at the medium
scale is small (2.5\%) and that we have not measured run-to-run
variance on the perplexity numbers themselves. The one downstream task
where we did report multi-seed numbers --- the Olist review probe of
Section~\ref{sec:olist} --- placed the two architectures within
$1\sigma$ of each other and inverted the IMDB-pt ranking, which
underlines that a definitive claim about ordering at this scale would
require multi-seed pretraining runs as well; we view our observation
as suggestive rather than conclusive.

\paragraph{Why the Transformer wins at the tiny scale.}
The original xLSTM paper~\cite{beck2024} reports advantages over
Transformers from $\sim$125\,M parameters and tens of billions of
training tokens upward, and never below this scale. Our tiny scale
($\sim$1.45\,M parameters, $\sim$1.5\,M tokens) lies three to four
orders of magnitude below any setting in which the architecture has
been validated. The xLSTM block carries machinery that the Transformer
does not---block-diagonal q/k/v projections, learnable skip connection,
component-wise output gates, exponential gating with stabiliser state,
GroupNorm head-wise---which costs parameters that must themselves be
learned from data. In the regime we study, where the token-to-parameter
ratio is one to two orders of magnitude below
Chinchilla-optimal~\cite{hoffmann2022}, the marginal cost of carrying
this gating apparatus is high relative to the marginal benefit. A
purely-attention-plus-MLP architecture concentrates more of its budget
in primitives that are simpler to learn from limited data. The tiny-scale
gap is therefore not anomalous; rather, it is the predicted outcome of
applying an architecture rich in inductive bias far outside its
data-and-capacity envelope.

\paragraph{What the medium-scale crossover supports.}
The crossover is in the direction predicted by the
scale-dependent-advantage hypothesis: as we move from tiny to medium,
the Transformer's $28\%$ lead inverts to a $2.5\%$ deficit. Even with
the seed-variance caveat, the \emph{direction} of the change is more
robust than the magnitude. What is non-trivial is that the inversion
occurs as early as $\sim$4.5\,M parameters, two orders of magnitude
below the smallest scale at which the original paper documents the
effect. We therefore interpret our medium-scale result not as a
challenge to Beck et al.\ but as a lower-bound case study extending the
observed envelope of the effect downward. We make no claim about the
asymptotic advantage at scale.

\paragraph{The GPU sweep splits the perplexity and downstream signals.}
The GPU sweep of Section~\ref{sec:gpu-scaling} extends the case study
to $10$\,M parameters and a larger training budget, and surfaces a
qualitatively new pattern: at $10$\,M, the Transformer regains the lead
on held-out perplexity ($269$ vs.\ $287$, a $6\%$ gap) yet \emph{loses}
on both downstream probes we measure (author classification by
$2.2$ accuracy points, IMDB-PT sentiment by $0.8$). The xLSTM $10$\,M
checkpoint also overfits the corpus much earlier than the Transformer
$10$\,M checkpoint (best step $1{,}600$ vs.\ $4{,}400$ of $8{,}000$),
which suggests the xLSTM is more data-hungry at this scale --- consistent
with the inductive-bias argument made for the tiny-scale result above,
now playing out in the opposite regime (capacity exceeding effective
corpus entropy rather than corpus exceeding learnable capacity). The
practical takeaway for a reader picking an architecture at this scale
is that perplexity alone does not arbitrate: a model with a lower next-token
loss may not produce more linearly-separable representations for
downstream tasks. The Olist multi-seed result of
Section~\ref{sec:olist} additionally warns that the IMDB-PT gap is
within the seed-to-seed band of a similar binary-sentiment task, so the
two architectures are best read as at statistical parity on out-of-domain
sentiment, with the xLSTM holding a thin lead.

\paragraph{Scope of the comparison.}
The matched-LR design (peak $3 \times 10^{-4}$ for both architectures)
prioritises symmetry over per-architecture tuning. The original xLSTM
uses higher peak LRs at small scale
($3\times10^{-3}$ at 125\,M~\cite{beck2024}); the present comparison is
therefore a matched-hyperparameter comparison rather than a
best-of-each-architecture comparison. The token-to-parameter ratio of
the medium runs ($\approx 0.45$) is well below the $\approx 20{:}1$
Chinchilla ratio, placing the comparison in the data-limited rather than
the capacity-limited regime; this is an intentional property of
the corpus (the public-domain 19th-century Portuguese literary record)
rather than an artefact of the training schedule. Section~\ref{sec:limitations}
makes the explicit scope of the reported numbers concrete.

% =============================================================================
\section{Limitations and Future Work}\label{sec:limitations}

\paragraph{Single seed per configuration.}
The CPU numbers of Table~\ref{tab:results} and the GPU numbers of
Table~\ref{tab:gpu-results} are each from one training run per
configuration. The single-digit-percent perplexity gaps at medium and
$10$\,M scales are within the seed-variance band typical of small-scale
language modelling. The Olist multi-seed downstream probe of
Section~\ref{sec:olist} ($\pm 1\%$ accuracy between seeds on the same
checkpoint) is the only measurement in this paper with empirical error
bars; the other tables should be read as point estimates and their
rankings as ordinal rather than as precise effect sizes.

\paragraph{Domain.}
The corpus is restricted to 19th-century Portuguese-language literature.
The reported observations are scoped to this regime; no claim is made
about transfer to modern Portuguese, technical text, code, or other
languages.

\paragraph{Holdout author overlap.}
The held-out work (\emph{Memórias Póstumas de Brás Cubas}) shares its
author with six of the nineteen training novels, Machado de Assis being
the most prolific author of the period in the public domain. The works
themselves are disjoint at the title level, but author-level disjointness
is not achieved.

% =============================================================================
\section{Conclusion}

We have presented an open-source pipeline that isolates architecture as
the sole variable in small-scale language modelling experiments. As an
illustrative case study, we have applied the pipeline to xLSTM and a
GPT-style Transformer at five parameter scales across two hardware
regimes, using a corpus of 19th-century Portuguese literature. On the
CPU sweep we observed a crossover in held-out perplexity --- the
Transformer wins at the smallest scale, the xLSTM marginally surpasses
the Transformer at the medium scale --- whose direction is consistent
with the scale-dependent-advantage hypothesis of Beck et
al.~\cite{beck2024}, with the additional empirical observation that
the inversion is already present at $\sim$$4.5$\,M parameters, two
orders of magnitude below the smallest scale at which the original
work documents the effect. The GPU sweep extends the picture: the
Transformer regains the lead in held-out perplexity at $10$\,M
parameters ($269$ vs.\ $287$, a $6\%$ gap), but the xLSTM overfits
markedly earlier than the Transformer in this regime, and on every
downstream probe we measured (author classification, IMDB-PT
sentiment, Olist review-score) the xLSTM matches or marginally
surpasses the Transformer at every scale considered. Taken together,
the cross-architecture ordering depends on the metric: perplexity
favours the Transformer at the largest scale, downstream probes favour
the xLSTM, and these two answers should be read together rather than
arbitrated. We additionally highlight a methodological subtlety that is
rarely controlled in published arch-comparison work: training
throughput differences between recurrent and attention-based
architectures in our setting span $\sim$$5$--$8\times$ on CPU and
$\sim$$45$--$80\times$ on GPU due to backend implementation rather than
fundamental architectural cost.

% =============================================================================
\section*{Code and Data Availability}

The full pipeline, including all configuration files, training
scripts, analysis scripts, and the trained checkpoints from this study,
is available at\\
\url{https://github.com/Joao-pedrosantos/NaturalLanguageProcessing}.
All training data is publicly available from Project Gutenberg under
public-domain or non-restrictive licenses.

% =============================================================================
\begin{thebibliography}{9}

\bibitem{vaswani2017}
Vaswani, A. \emph{et al.}
Attention is all you need.
\emph{Advances in Neural Information Processing Systems 30} (2017).

\bibitem{devlin2019}
Devlin, J., Chang, M.-W., Lee, K. \& Toutanova, K.
BERT: Pre-training of deep bidirectional transformers for language
understanding.
\emph{NAACL-HLT} (2019).

\bibitem{beck2024}
Beck, M. \emph{et al.}
xLSTM: Extended long short-term memory.
\emph{Advances in Neural Information Processing Systems 37} (2024).

\bibitem{radford2019}
Radford, A. \emph{et al.}
Language models are unsupervised multitask learners.
OpenAI technical report (2019).

\bibitem{sennrich2016}
Sennrich, R., Haddow, B. \& Birch, A.
Neural machine translation of rare words with subword units.
\emph{Proceedings of ACL} (2016).

\bibitem{hoffmann2022}
Hoffmann, J. \emph{et al.}
Training compute-optimal large language models.
\emph{Advances in Neural Information Processing Systems 35} (2022).

\bibitem{nanogpt}
Karpathy, A.
nanoGPT: The simplest, fastest repository for training/finetuning
medium-sized GPTs.
GitHub repository (2022). \url{https://github.com/karpathy/nanoGPT}.

\bibitem{maas2011imdb}
Maas, A.~L. \emph{et al.}
Learning word vectors for sentiment analysis.
\emph{Proceedings of ACL-HLT} (2011).

\bibitem{huggingfacetokenizers}
Hugging Face.
\texttt{tokenizers}: Fast state-of-the-art tokenizers optimised for
research and production.
GitHub repository.
\url{https://github.com/huggingface/tokenizers}.

\bibitem{tiedemann2020opus}
Tiedemann, J. \& Thottingal, S.
OPUS-MT --- building open translation services for the world.
\emph{Proceedings of the 22nd Annual Conference of the European
Association for Machine Translation (EAMT)} (2020).

\bibitem{junczys2018marian}
Junczys-Dowmunt, M. \emph{et al.}
Marian: Fast neural machine translation in C++.
\emph{Proceedings of ACL 2018, System Demonstrations} (2018).

\bibitem{dao2022flashattention}
Dao, T., Fu, D.~Y., Ermon, S., Rudra, A. \& Ré, C.
FlashAttention: Fast and memory-efficient exact attention with
IO-awareness.
\emph{Advances in Neural Information Processing Systems 35} (2022).

\end{thebibliography}

% =============================================================================
\appendix
\section{Generated Samples}\label{app:samples}

All samples are generated with temperature 0.8 and top-$k$ 40, conditioned on
the prompt \texttt{"Capitu olhou para mim"}. We reproduce the model output
verbatim; line breaks reflect the model's own newline tokens.

\paragraph{xLSTM tiny (1.46\,M).}
\begin{quote}\small
\itshape
Capitu olhou para mim. E ``I?'' Eu, ``S and I said, ``All to make
--Não, com a to and we had, que me as, e a ser em lir a Cas. --V O seus
ret knowra, e os meus bo-se, com ella, a cabeça, a cim, mas, n'um que é
que o pae de Capitú ao Ca,
\end{quote}

\paragraph{Transformer tiny (1.45\,M).}
\begin{quote}\small
\itshape
Capitu olhou para mim, peu á Europa; elle me fez um de uma fóimente o
nosso velho, não teria a primeira, e Conversinaria-se á vida, que eu ia
do bureira. O APITULO No dia, de CA E a mãe, e o que ella Hou-se, quando
se os lugará
\end{quote}

\paragraph{xLSTM medium (4.55\,M).}
\begin{quote}\small
\itshape
Capitu olhou para mim a passel. Os que eu não ser não a minha mãe, a
minha mãe, em dirá á porta de dizer. Ora, mas não é preciso que eram que
não lhe e eu mesmo. Vinha-me com a verdade com seu filho, e não ha de
profeira. --Ha? Não, não ha de mim. --Aqui, mas não se não ha do outro
de Deus. --Não,
\end{quote}

\paragraph{Transformer medium (4.22\,M).}
\begin{quote}\small
\itshape
Capitu olhou para mim, peu á Europa; elle me fez um de uma fóimente o
nosso velho, não teria a primeira, e Conversinaria-se á vida, que eu
ia do bureira. O APITULO No dia, de CA E a mãe, e o que ella Hou-se,
quando se os lugará
\end{quote}

\medskip

We note qualitatively that medium-scale samples reproduce 19th-century
Portuguese orthographic conventions (\emph{pae}, \emph{á}, \emph{ella},
contractions like \emph{n'um}) and recognisable dialogue structure with
em-dashes. Tiny-scale samples produce isolated correct words but lower
syntactic coherence. We refrain from quantitative qualitative claims.

\subsection*{Multi-prompt completions (medium scale)}\label{app:multiprompt}

To probe behaviour beyond a single conditioning context, we additionally
sampled both medium-scale models on four distinct prompts spanning
narrative openings and dialogue turns. All completions are generated with
the same decoding settings as above (temperature 0.8, top-$k$ 40,
\texttt{torch.manual\_seed(42)} reset before each prompt) and 80 new
tokens. Outputs are reproduced verbatim; line breaks are the model's own.

\paragraph{xLSTM medium (4.55\,M).}

\noindent\emph{Prompt: \texttt{"Capitu olhou para mim"}}
\begin{quote}\small\itshape
Capitu olhou para mim a mulher, e a minha mãe, ao que o que não me
dissesse e não era a mãe.\\
--Java é que ha ser o que a voz.\\
--Não é que me parece.\\
--É é verdade.\\
--Nem eu?\\
--Fou não me deu ao marido.
\end{quote}

\noindent\emph{Prompt: \texttt{"O sol nascia sobre o"}}
\begin{quote}\small\itshape
O sol nascia sobre o palavras.\\
No dia, que o que não me não sei de Sativa.\\
--Sim não? Mas é que ha ser o que a voz de minha mãe me palavras.\\
--Não tambem é verdade.\\
--Nem eu?\\
--Fou não me deu ao marido.
\end{quote}

\noindent\emph{Prompt: \texttt{"Era uma vez, em uma cidade pequena,"}}
\begin{quote}\small\itshape
Era uma vez, em uma cidade pequena, e não se lhe não tinha ser que não
é que se não me dissesse e não era a mãe.\\
--Java é que ha ser o que a voz.\\
--Não é que me parece.\\
--É é verdade.
\end{quote}

\noindent\emph{Prompt: \texttt{"-{}-Bom dia, senhor."}}
\begin{quote}\small\itshape
--Bom dia, senhor.\\
--Saro mesmo, disse-me o que não me dissesse, disse ella, mas ha para
o corpo. Mas é que se lhe lembra que a voz debas accando. O outro é
verdade, porque é que me não não posso, eu o a cabeça do ceu, a
palavra de mim.
\end{quote}

\paragraph{Transformer medium (4.22\,M).}

\noindent\emph{Prompt: \texttt{"Capitu olhou para mim"}}
\begin{quote}\small\itshape
Capitu olhou para mim é me que eu é preciso ser que não é que se não me
duas. Um dia? Mas é que se ser o que a voz de CAlexarminha ao que a
verdade.\\
--Não me deu ao marido.\\
--Como.
\end{quote}

\noindent\emph{Prompt: \texttt{"O sol nascia sobre o"}}
\begin{quote}\small\itshape
O sol nascia sobre o ezer-o, como a casa que o que não me dão. Um dia?
Mas é que se ser o que a voz de CAjam o que ella, não é verdade.\\
--Não me deu ao marido. Os labios, não.
\end{quote}

\noindent\emph{Prompt: \texttt{"Era uma vez, em uma cidade pequena,"}}
\begin{quote}\small\itshape
Era uma vez, em uma cidade pequena, eja não lhe é preciso ser que não
que o que não me postou, quando mas é para o que não podia ser tudo.
que eu a voz e de a vez que é que me o outro é verdade.
\end{quote}

\noindent\emph{Prompt: \texttt{"-{}-Bom dia, senhor."}}
\begin{quote}\small\itshape
--Bom dia, senhor.\\
--CAPa, mas eu não é, quando e para o que não não é que se ser o que
a voz e de preço, e não se o que se tambem é verdade.\\
--Que?\\
--Não me deu ao marido.\\
--Pois me não.
\end{quote}

\medskip

These completions are consistent with the held-out perplexity numbers in
Table~\ref{tab:results}: both models stay on-distribution lexically but
collapse syntactically after a few clauses, and both repeat short
fragments (\emph{``Não me deu ao marido''} appears across prompts in
both architectures). The Transformer's outputs contain more whitespace
artefacts and more aggressively segmented dialogue, while the xLSTM's
outputs flow into longer pseudo-sentences. We stress that, as throughout
this section, we draw no quantitative conclusion from these samples;
they are included to make the qualitative regime of the medium-scale
runs concrete for the reader.

\end{document}
